\begin{papercard}{feinstein1990}
\cf{Stages} \stagetag{I}. Part I of the classic pair on the ``paradoxes of kappa'': an algebraic and worked-table demonstration that $\kappa$ can be low despite high observed agreement, and higher for \emph{worse}-looking (asymmetric) marginals. Part II (Cicchetti \& Feinstein, not in this corpus) proposes remedies. The authors' position is critical: they would rather drop the chance correction and design a balanced challenge population.

\cf{Setting} Binary ratings by two observers in a fourfold table with cells $a,b,c,d$ and margins $f_1,f_2$ (A) and $g_1,g_2$ (B). Observed agreement $p_o=(a+d)/N$; expected $p_e=(f_1g_1+f_2g_2)/N^2$; $\kappa=(p_o-p_e)/(1-p_e)$ \pc{feinstein1990}{2}{For the calculation of K, a ratio is formed between the chance-corrected observed agreement}. Like predictive values in diagnostic testing, $\kappa$ is affected by prevalence, in the form of the marginal totals (Kraemer, Grove et al.). The authors consider only the reliability use of $\kappa$, where no gold standard exists \pc{feinstein1990}{2}{we are concerned exclusively with the much more common use of K as an index of “reliability”}.

\cf{Paradox 1: high $p_o$, low $\kappa$} For $p_o=.85$: with balanced margins ($p_e=.50$) $\kappa=.70$; with imbalanced margins ($p_e=.78$) $\kappa=.32$. Tables 1 vs.\ 2 have identical agreement (40+45 vs.\ 80+5 of 100) but $\kappa=.70$ vs.\ $.32$ \pc{feinstein1990}{2}{if pe is large, the correction process can convert a relatively high value of p. into a relatively low value of K}. $p_e=vw+(1-v)(1-w)$ with each observer's ``prevalence'' $v,w$: with both far from 0.5 on the same side (symmetrical imbalance) $p_e=0.5+2v'w'$ grows, so $\kappa$ falls \pc{feinstein1990}{3}{Thus, in the first circumstance, with symmetrical imbalances in the marginal totals}.

\cf{Paradox 2: asymmetry is rewarded} If observers' margins diverge in opposite directions, $p_e=0.5-2v'w'$ shrinks and $\kappa$ rises. Table 3 (margins 70/30 vs.\ 60/40) has $\kappa=.13$, while Table 4 (30/70 vs.\ 60/40), apparently worse, has $\kappa=.26$ \pc{feinstein1990}{3}{Table 4, despite its apparently worse marginal-total agreement, thus emerges with a higher value for K than Table 3}. Tables 5 and 6: both $p_o=.90$ with symmetrically unbalanced margins, but perfect symmetry gives $\kappa=.44$ and imperfect symmetry gives $.74$ \pc{feinstein1990}{4}{In both cases of unbalanced relatively symmetrical marginals, the impressive p. of 0.90 is substantially lowered by K}. The alternative algebra $\kappa=\dfrac{N(a+d)-2f_1f_2-(f_1-f_2)(a-d)}{N^2-2f_1f_2-(f_1-f_2)(a-d)}$ shows that the sign of $(f_1-f_2)(a-d)$ decides: concordant increments reduce $\kappa$, discordant ones raise it (Appendix).

\cf{Marginals bound the maximum} The maximum possible $p_o$ is (smallest margin + its diagonally opposite margin)$/N$. High $p_o$ is only possible when imbalances are (nearly) symmetric; strongly asymmetric imbalances cap $p_o$ (Table 4: .70) \pc{feinstein1990}{4}{the maximum possible value for p. is (l/N) times the}. $\kappa_{\max}=(\max p_o-p_e)/(1-p_e)$ was considered as a repair, but it ``does not repair either difficulty, and seems to make the second one worse'': it leaves paradox 1 and amplifies paradox 2, raising Table 4's $\kappa$ from .26 to .58 and Table 6's from .74 to 1.00 \pc{feinstein1990}{6}{An adjustment that substitutes K max for K does not repair either difficulty}.

\cf{Critique of the chance correction} The correction assumes each observer has a relatively fixed prior probability of responding yes or no, which seems inappropriate for most clinical observers who respond to whatever is presented, unless they know the population is mostly normal or abnormal. Since the correction depends on the marginals, observers can get low $\kappa$ at high $p_o$ and are ``penalised'' for a lopsided challenge population (Tables 2, 5) and ``rewarded'' for opposite-direction imbalances (Tables 4, 6) \pc{feinstein1990}{6}{The presumptive reason for the chance correction in K is that the measuring instruments are human observers}. Proposal: drop $\kappa$, score observers by $p_o$, and make the investigator choose a balanced population (a logistic burden many would reject); Part II offers alternative indexes and a strategy \pc{feinstein1990}{6}{The removal of the K coefficient would eliminate these “unfair” corrections for chance}.

\cf{Use for this thesis} The definitive description of why $\kappa$ can mislead when a code is rare or ubiquitous, which is typical for dark-pattern or consent-quality codes. Operational advice: (1) always publish the 2$\times$2 (or $k\times k$) tables and the marginals with any $\kappa$; (2) report raw agreement and positive/negative specific agreement alongside; (3) design the pilot sample so each code occurs often enough (purposive inclusion of apps with and without the feature); (4) decide beforehand how a high-$p_o$/low-$\kappa$ result will be handled. The authors' remedy (abandoning chance correction) is contested; Krippendorff argues the opposite (the data lack variation, so low coefficients are correct) and Gwet proposes AC1, so the thesis must pick and justify its position and cite Part II (to be added) for the formal remedies.
\end{papercard}
